\chapter{Conclusion}
\label{ch:conclusion}

We set out to evaluate Kuppa and Le-Khac's Contrastive NCM drift detector~\cite{Kuppa_2022} not in isolation, as the literature does, but coupled to the downstream base learner it is meant to supervise -- the configuration in which it would actually be deployed. That coupling, not the detector's representation, is where the security of the system is decided.

Our central finding is that the coupling exposes the adversary to the system's adaptation rather than to its learned decision boundary directly. That is, whether and when the base learner is rebuilt, and what is added to it at a rebuild. Three results make the case. First (RQ1), because the base learner changes only at discrete concept-discovery events, an adversary who controls the firing predicate, or the buffer at a fire, controls what the learner learns and if it learns at all; a continuously-updated learner offers no such gate. Second (RQ3), the firing trigger is \emph{forceable and not robust}: an aimed feature-space mean-shift forces a concept-discovery event both white-box and grey-box, retaining roughly $70\%$ of its displacement against a data-disjoint surrogate, while the per-sample detection it is built on transfers poorly. Thus, it is the batch-mean statistic the detector adopts for stability that makes the trigger transferable. The firing floors that resist forcing are themselves double-edged: the gate ablation proves that any floor on the buffer's makeup which blocks forcing simultaneously enables suppression, an over-determined trade-off intrinsic to makeup-based gating. Third (RQ2), the composition channel resists poisoning that contradicts $\mathcal{M}$'s decision boundary: a from-scratch learner holds its \texttt{DoS} false-negative rate near its clean $0.02$ across injection budgets. However, it does \emph{not} resist a backdoor: a benign-labelled, trigger-stamped payload installs a shortcut that evades $\mathcal{M}$ on demand (a false-negative rate of $0.69$ on triggered traffic against $0.03$ on clean), leaves aggregate accuracy intact, and persists via the never-evict exemplar memory. We also find that the same channel governs whether the learner adapts at all. A benign dilution of $p = 25\%$ holds the gate shut for the whole stream and collapses the system to the no-adaptation floor, $F_1$ falling from $0.71$ to $0.207$. Finally, the system never discovers the second novel class -- but not because the detector fails to flag it. \texttt{PortScan} clears both the count and displacement conditions; it is blocked only by our $0.30$ fraction floor, diluted below threshold by the detector's own benign false positives. The blind spot is thus the suppression mechanism arising with no adversary at all, an artefact of our gate's calibration rather than of the detector.

A key takeaway is that coupling a robust detector to a base learner does not guarantee that the system will inherit the detector's robustness; it creates new attack surfaces that the detector-only evaluation the original proposal invites cannot see. A drift detector that is robust as a classifier can still be an exploitable \emph{trigger} or \emph{supervisor}.

Several broader lessons follow. The first is that \emph{robustness does not compose}: a component certified robust in isolation can be redeployed as a trigger or a supervisor whose failure modes its own evaluation never measured, so security has to be assessed at the level of the coupled system rather than the component. The second is that the design choices made for good behaviour are exactly what the adversary turns against the system -- the batch-mean statistic adopted to make the trigger \emph{stable} is what makes it transferable, and the rehearsal memory adopted to prevent catastrophic \emph{forgetting} is what makes a backdoor durable. The third is that a defended decision boundary is not a defended system: with the boundary held, the adversary's leverage simply moves to whether and exactly when the learner is rebuilt, and to whichever class the detector never learns to flag. Finally, the experiments are a reminder that aggregate metrics conceal precisely the damage that matters here; per-class reporting was necessary to see a whole attack class go undetected behind an almost-unchanged $F_1$. Each is a factor that designers of coupled adaptive systems must weigh, and each is invisible to the detector-only evaluation the current research relies on.

\section{Future Work}
\label{sec:conc-future}
Several threads follow directly. A \emph{feature-constrained} mean-shift, restricting perturbations to \emph{protocol-valid} flows, would establish how forceable the trigger is for a realistic attacker, since our attack succeeds only at an unconstrained $L_\infty$ budget. The backdoor we demonstrate uses a synthetic feature trigger; realising a protocol-valid trigger on genuine flows and defending against it, for instance by sanitising the harvested exemplars before they are replayed, is the natural follow-up. The \emph{blind spot} invites a \emph{stealth-novelty} attack: since a novel class is admitted only once its flows make up enough of the drifted buffer, an adversary could pace a novel attack under benign cover so that its attack fraction never reaches the gate's fraction floor, keeping the coupled system from ever adapting to it. Finally, from a defensive perspective, resetting or windowing the displacement accumulator on gated crossings would likely close the count-floor suppression mechanism, and a continuously-updated base-learner baseline would quantify what the discrete trigger costs in robustness.

\section{Limitations}
\label{sec:conc-limitations}
The findings of~\cref{ch:evaluation} are reported at a single calibrated operating point. The per-sample drift threshold $T_{drift}$ is fixed to a $10\%$ known-class false-positive rate and the concept threshold $T$ to the $0.995$ displacement quantile (\cref{sec:impl-calibration}); both are principled choices, but the attack results are necessarily conditional on them. Because Kuppa and Le-Khac underspecify their own operating point, we could not validate these thresholds by reproduction (which is why the detector itself is validated threshold-free, on AUROC, in~\cref{sec:impl-repro-results}); we validate them instead by the system's correct adaptation under genuine drift. A systematic sensitivity analysis -- sweeping $T_{drift}$ across false-positive rates and $T$ across quantiles -- would establish how far the findings depend on this calibration, and is left to future work.

The firing floors (count $500$, fraction $0.3$) are likewise calibrated rather than optimised: the fraction floor is set above the detector's $\sim$$10\%$ false-positive contamination of the buffer and the count floor at a minimum viable sample for a stable prototype (\cref{sec:impl-gate}). Our claims are scoped accordingly. The over-determination of the makeup gate is argued structurally and so does not depend on these values (\cref{sec:eval-rq3-gate}), whereas the magnitude of each attack's budget, and the \texttt{PortScan} blind spot, are conditional on this operating point. A sweep of the floors would map the force-versus-suppress trade-off curve. However, because our contribution is to characterise the attack surface rather than to optimise a defence, we leave that, and the choice of an operating point that best balances the two, to future work.

